\subsection{BOUNDS OF A REAL-VALUED FUNCTION}

DEFINITION 3. Let \(f\) be a real-valued function defined on a set \(\mathbf{X}\) , and let \(\mathbf{A}\) be a non-empty subset of \(\mathbf{X}\) . Then the least upper bound (resp. greatest lower bound) of the set \(f(\mathbf{A})\) in \(\overline{\mathbf{R}}\) is called the least upper bound (resp. greatest lower bound) of \(f\) in \(\mathbf{A}\) , and is denoted by \(\sup_{x\in \mathbf{A}}f(x)\left[\text{resp. } \inf_{x\in \mathbf{A}} f(x)\right]\) .

In particular, if A is a non-empty subset of \(\overline{R}\) , then

\[
\sup \mathrm{A} = \sup _ {x \in \mathrm{A}} x.\tag{1}
\]

It is often more convenient to use the notation on the right-hand side to denote the least upper bound of A.

The number \(a = \sup_{x \in \mathbf{A}} f(x)\) is characterized by the following two properties:

\begin{enumerate}
\def\labelenumi{(\roman{enumi})}
\item
  For all \(x \in \mathbf{A}, f(x) \leqslant a\) .
\item
  For every \(b < a\) , there exists \(x \in \mathbf{A}\) such that \(b < f(x) \leqslant a\) .
\end{enumerate}

The numbers \(\sup_{x\in\mathbf{A}}f(x)\) and \(\inf_{x\in\mathbf{A}}f(x)\) belong to the closure of \(f(\mathbf{A})\) in \(\overline{\mathbf{R}}\) . We have \(\inf_{x\in\mathbf{A}}f(x)\leqslant\sup_{x\in\mathbf{A}}f(x)\) ; and these two numbers are equal if and only if \(f\) is constant on \(\mathbf{A}\) .

A real-valued function, defined on a set X, is bounded above (resp. bounded below) in a non-empty subset A of X if and only if \(\sup_{x\in\mathbf{A}}f(x)<+\infty\) {[}resp. \(\inf_{x\in\mathbf{A}}f(x)>-\infty\) {]}. f is bounded in A if and only if \(|f|\) is bounded above in A, hence if and only if \(\sup_{x\in\mathbf{A}}|f(x)|<+\infty\) .

We have

\[
\inf _ {x \in \mathbf {A}} f (x) = - \sup _ {x \in \mathbf {A}} (- f (x)).\tag{2}
\]

This relation reduces all properties of the greatest lower bound to those of the least upper bound; hence in general we shall speak only of the latter.

PROPOSITION 5. Let \(f\) be a real-valued function defined on a set \(\mathbf{X}\) . On the set \(\mathfrak{F}(\mathbf{X})\) of all finite subsets of \(\mathbf{X}\) , directed with respect to the relation \(\subset\) , the real-valued function \(\mathbf{H} \to \sup_{x \in \mathbf{H}} f(x)\) is increasing, the real-valued function \(\mathbf{H} \to \inf_{x \in \mathbf{H}} f(x)\) is decreasing, and we have

\[
\left\{ \begin{array}{l} \sup _ {x \in \mathbf {A}} f (x) = \lim _ {\mathbf {H} \in \mathfrak {F} (\mathbf {X})} (\sup _ {x \in \mathbf {H}} f (x)), \\ \inf _ {x \in \mathbf {A}} f (x) = \lim _ {\mathbf {H} \in \mathfrak {F} (\mathbf {X})} (\inf _ {x \in \mathbf {H}} f (x)). \end{array} \right.\tag{3}
\]

Let \(\varphi(\mathbf{H}) = \sup_{x \in \mathbf{H}} f(x)\) . Clearly \(\varphi\) is increasing, and therefore has a limit \(a\) (no. 2, Theorem 2); and since \(\varphi(\mathbf{H}) \leqslant \sup_{x \in \mathbf{A}} f(x)\) for all H, we have \(a \leqslant \sup_{x \in \mathbf{A}} f(x)\) (no. 2, Theorem 1). If we had \(a < \sup_{x \in \mathbf{A}} f(x)\) , then there would exist \(x_0 \in \mathbf{X}\) such that \(a < f(x_0)\) ; but this is a contradiction since \(\varphi(\mathbf{H}) \geqslant f(x_0)\) whenever \(x_0 \in \mathbf{H}\) .

In particular, by (1), if A is any non-empty subset of \(\overline{R}\) , we have

\[
\sup \mathrm{A} = \lim _ {\mathrm{H} \in \mathfrak {F} (\mathrm{A})} (\sup _ {x \in \mathrm{H}} x).\tag{4}
\]

PROPOSITION 6. Let \(f\) and \(g\) be two real-valued functions defined on \(\mathbf{X}\) . If \(f(x) \leqslant g(x)\) at every point \(x\) of a non-empty subset \(A\) of \(\mathbf{X}\) , then we have

\[
\left\{ \begin{array}{l} \sup _ {x \in \mathbf {A}} f (x) \leqslant \sup _ {x \in \mathbf {A}} g (x), \\ \inf _ {x \in \mathbf {A}} f (x) \leqslant \inf _ {x \in \mathbf {A}} g (x). \end{array} \right.\tag{5}
\]

PROPOSITION 7. Let \(f\) be a real-valued function defined on \(\mathbf{X}\) . If \(\mathbf{A}\) and \(\mathbf{B}\) are two non-empty subsets of \(\mathbf{X}\) such that \(\mathbf{A} \subset \mathbf{B}\) , then

\[
\sup _ {x \in \mathbf {A}} f (x) \leqslant \sup _ {x \in \mathbf {B}} f (x).\tag{6}
\]

PROPOSITION 8. Let \(f\) be a real-valued function defined on \(\mathbf{X}\) , and let \((\mathbf{A}_t)_{t \in I}\) be a non-empty family of non-empty subsets of \(\mathbf{X}\) ; then

\[
\sup_{x\in \bigcup_{t\in \mathbf{I}}\mathbf{A}_{t}}f(x) = \sup_{t\in \mathbf{I}}(\sup_{x\in \mathbf{A}_{t}}f(x)).\tag{7}
\]

Let \(f\) be a real-valued function defined on a product set \(\mathbf{X}_1 \times \mathbf{X}_2\) . If \(\mathbf{A}_2\) is a non-empty subset of \(\mathbf{X}_2\) we shall denote by \(\sup_{x_2 \in \mathbf{A}_2} f(x_1, x_2)\) the least upper bound in \(\mathbf{A}_2\) of the real-valued function \(x_2 \to f(x_1, x_2)\) defined on \(\mathbf{X}_2\) . From Proposition 8 we deduce in particular:

PROPOSITION 9. Let \(f\) be a real-valued function defined on a product set \(\mathbf{X}_1 \times \mathbf{X}_2\) . If \(\mathbf{A}_1, \mathbf{A}_2\) are any non-empty subsets of \(\mathbf{X}_1, \mathbf{X}_2\) respectively, then

\[
\sup _ {(x _ {1}, x _ {2}) \in \mathbf {A} _ {1} \times \mathbf {A} _ {2}} f (x _ {1}, x _ {2}) = \sup _ {x _ {1} \in \mathbf {A} _ {1}} (\sup _ {x _ {2} \in \mathbf {A} _ {2}} f (x _ {1}, x _ {2})) = \sup _ {x _ {2} \in \mathbf {A} _ {2}} (\sup _ {x _ {1} \in \mathbf {A} _ {1}} f (x _ {1}, x _ {2})).\tag{8}
\]

